\label{chap:pt2_state_of_the_art}

The use of neural networks for solving physical problems has been extensively studied in the literature, particularly following the introduction of Physics-Informed Neural Networks (PINNs) by Raissi et al.~\cite{raissi2019physics}. 

A significant portion of PINN approaches and data-driven methods in thermoforming has focused on inverse modeling strategies aimed at identifying optimal process parameters to reduce defects in the final product geometry. For instance, Yang et al.~\cite{yang2004modeling} proposed an inverse neural network model to predict optimal process conditions for achieving a uniform thickness distribution across the sheet. Chang et al.~\cite{chang2005derivation} developed an inverse model that maps desired product dimensions to process parameters in order to minimize dimensional deviations. Similarly, Leite et al.~\cite{leite2018vacuum} trained a neural network using processing parameters as inputs and defect indicators as outputs, incorporating an optimization objective to reduce forming defects. Tan et al.~\cite{tan2022prediction} employed dual neural networks to infer process parameters from simulation images and optimize them for shear angle uniformity.

More recent studies have shifted attention toward earlier stages of the thermoforming process, particularly the heating phase. Simoncini et al.~\cite{simoncini2017neural} proposed a neural network model to predict polymer behavior under coupled thermal and mechanical effects, with emphasis on infrared heating interactions. Römer et al.~\cite{romer2022temperature} applied reinforcement learning to design a temperature control strategy for the pre-forming stage in automated tape laying processes. Jalilvand et al.~\cite{jalilvand2024towards} used deep reinforcement learning to optimize heating settings for a target temperature distribution, achieving reduced energy consumption. Ayçiçek et al.~\cite{aycciccek2024ai} developed an artificial intelligence-based thermal inspection system using thermal imaging, reducing defects and material waste. Turan et al.~\cite{turan2025convolutional} proposed a surrogate model based on Convolutional Neural Networks (CNNs) and addressed the inverse problem of predicting lamp powers using gradient-based optimization.

Despite these advances, existing approaches either rely on surface temperature information or do not explicitly model the temperature distribution at the core of the sheet, which is critical for product quality. Moreover, these methods often assume simplified industrial conditions, such as uniform sheet properties or single-sided heating configurations, and frequently rely on steady-state assumptions that fail to capture the transient dynamics of the heating phase, limiting their applicability in real industrial environments.

In this study, we address these limitations by explicitly modeling the temperature distribution at the core of the sheet under realistic industrial thermoforming conditions. In contrast to prior work, the proposed method does not rely on steady-state assumptions; instead, it is trained on time-dependent simulations to accurately predict the final temperature distribution.